\begin{center}
\hypertarget{native-051-article033}{}\pdfbookmark[2]{33.}{article-033}33.
\end{center}
In iis igitur sectionibus conicis, quarum excentricitas ab unitate parum differt, i. e. in ellipsibus et hyperbolis, quae ad parabolam proxime accedunt, methodi supra expositae tum pro determinatione anomaliae verae e tempore, tum pro determinatione temporis ex anomalia vera\footnote{Quoniam tempus implicat factorem $a^{\frac{3}{2}}$ vel $b^{\frac{3}{2}}$, error in $M$ vel $N$ commissus eo magis augetur, quo maior fuerit $a=\frac{p}{1-ee}$, vel $b=\frac{p}{ee-1}$\raisebox{.45ex}{.}}, omnem quae desiderari posset praecisionem non patiuntur: quin adeo errores inevitabiles, crescentes dum orbita magis ad parabolae similitudinem vergit, tandem omnes limites egrederentur. Tabulae maiores ad plures quam septem figuras constructae hanc incertitudinem diminuerent quidem, sed non tollerent, nec impedirent, quominus omnes limites superaret, simulac orbita ad parabolam nimis prope accederet. Praeterea methodi supra traditae in hocce casu satis molestae fiunt, quoniam pars earum indirecta tentamina saepius repetita requirit: cuius incommodi taedium vel gravius est, si tabulis maioribus operamur. Haud sane igitur superfluum erit, methodum peculiarem excolere, per quam in hoc casu incertitudinem illam evitare, soloque tabularum vulgarium adminiculo praecisionem sufficientem assequi liceat.

\begin{center}
\hypertarget{native-051-article034}{}\pdfbookmark[2]{34.}{article-034}34.
\end{center}
Methodus vulgaris, per quam istis incommodis remedium afferri solet, sequentibus principiis innititur. Respondeat in ellipsi vel hyperbola, cuius excentricitas $e$, semiparameter $p$ adeoque distantia in perihelio $=\frac{p}{1+e}=q$, \SourceLeaf{052}tempori post perihelium $t$ anomalia vera $v$; respondeat porro eidem tempori in parabola, cuius semiparameter $=2q$, sive distantia in perihelio $=q$, anomalia vera $w$, massa $\mu$ vel utrimque neglecta vel utrimque aequali supposita. Tunc patet haberi
\[
\int\frac{pp\,\mathrm{d}v}{(1+e\cos v)^2}:\int\frac{4qq\,\mathrm{d}w}{(1+\cos w)^2}=\BareRoot{p}:\BareRoot{2q}
\]
integralibus a $v=0$ et $w=0$ incipientibus, sive
\[
\int\frac{(1+e)^{\frac{3}{2}}\mathrm{d}v}{(1+e\cos v)^2\BareRoot{2}}=\int\frac{2\mathrm{d}w}{(1+\cos w)^2}\raisebox{.45ex}{.}
\]
Designando $\frac{1-e}{1+e}$ per $\alpha$, $\operatorname{tang}\frac{1}{2}v$ per $\theta$, integrale prius invenitur
\[
\begin{aligned}
={}&\BareRoot{(1+\alpha)}\cdot\{\theta+\tfrac{1}{3}\theta^3(1-2\alpha)-\tfrac{1}{5}\theta^5(2\alpha-3\alpha\alpha)\\
&\qquad+\tfrac{1}{7}\theta^7(3\alpha\alpha-4\alpha^3)-\text{etc.}\},
\end{aligned}
\]
posterius $=\operatorname{tang}\frac{1}{2}w+\frac{1}{3}\operatorname{tang}\frac{1}{2}w^3$. Ex hac aequatione facile est determinare $w$ per $\alpha$ et $v$, atque $v$ per $\alpha$ et $w$, adiumento serierum infinitarum: pro $\alpha$ si magis placet introduci potest $1-e=\frac{2\alpha}{1+\alpha}=\delta$. Quum manifesto pro $\alpha=0$ vel $\delta=0$ fiat $v=w$, hae series sequentem formam habebunt:
\[
\begin{aligned}
w&=v+\delta v^{\prime}+\delta\delta v^{\prime\prime}+\delta^3v^{\prime\prime\prime}+\text{etc.}\\
v&=w+\delta w^{\prime}+\delta\delta w^{\prime\prime}+\delta^3w^{\prime\prime\prime}+\text{etc.},
\end{aligned}
\]
ubi $v^{\prime}$, $v^{\prime\prime}$, $v^{\prime\prime\prime}$ etc. erunt functiones ipsius $v$, atque $w^{\prime}$, $w^{\prime\prime}$, $w^{\prime\prime\prime}$ etc. functiones ipsius $w$. Quoties $\delta$ est quantitas perparva, hae series celeriter convergent, paucique termini sufficient ad determinandum $w$ ex $v$, vel $v$ ex $w$. Ex $w$ invenitur $t$, vel $w$ ex $t$ eo quem supra pro motu parabolico explicavimus modo.
